\documentclass[11pt]{article}

\usepackage[utf8]{inputenc}
\usepackage[margin=1in]{geometry}
\usepackage{amsmath,amssymb}
\usepackage{booktabs}
\usepackage{graphicx}
\usepackage{listings}
\usepackage{xcolor}
\usepackage{hyperref}

\setlength{\parskip}{0.8em}
\setlength{\parindent}{0pt}

\title{Mini Project 2: Tight-Binding Model -- Daily Note Draft}
\author{Nathan Smith \& Reynaldo Salvador Mares Gil}
\date{\today\\[4pt]
\small Department of Physics and Astronomy, West Virginia University, Morgantown, WV, USA}

\begin{document}
\maketitle

\section*{Problem Statement}

The electronic structures of two (imaginary) two-dimensional crystals, A and B, are described by a tight-binding model with a two-orbital basis. The dependence of the overlap and hopping parameters, $\eta_{nm}$ and $t_{nm}$, on the distance $R$ between lattice sites is given by the following functions:

\begin{align*}
\eta_{11}(R) &= \left(1 + \frac{R}{d_0} + \frac{R^2}{3d_0^2}\right) e^{-\frac{R}{d_0}} \\[6pt]
\eta_{12}(R) &= \frac{1}{2}\left(1 + \frac{R}{d_0} + \frac{R^2}{64d_0^2}\right) e^{-\frac{R}{2d_0}} - \frac{1}{2}e^{-\frac{R}{d_0}} \\[6pt]
\eta_{22}(R) &= \left(1 + \frac{R}{2d_0} + \frac{R^2}{24d_0^2}\right) e^{-\frac{R}{2d_0}} \\[6pt]
t_{11}(R) &= -t_0 \left(1 + \frac{2R}{d_0}\right) e^{-\frac{R}{d_0}} \\[6pt]
t_{12}(R) &= -t_0 \left(1 + \frac{2R}{d_0}\right) e^{-\frac{R}{d_0}} + t_0 \left(1 + \frac{R}{16d_0}\right) e^{-\frac{R}{2d_0}} \\[6pt]
t_{22}(R) &= -t_0 \left(1 + \frac{R}{2d_0} - \frac{R^2}{24d_0^2}\right) e^{-\frac{R}{2d_0}}
\end{align*}

where $d_0$ and $t_0$ are positive constants that we use as length and energy units, respectively. The expressions for the hopping parameters are valid for $R > 3d_0$ and we have $t_{nm}(0) = 0$. The atomic energies associated with the two orbitals are $\epsilon_1 = 0$ and $\epsilon_2 = 0.8t_0$.

\begin{enumerate}
    \item Calculate the electronic band structures of the two crystals assuming that the lattice constant is: (i) $a = 5d_0$; (ii) $a = 7d_0$; (iii) $a = 9d_0$. Include the nearest neighbor and the next nearest neighbor contributions.
    \item Compare the results in (1) with the corresponding energy spectra calculated: (i) using only nearest neighbor contributions; (ii) neglecting the overlap integrals, $\eta_{nm}(R) = \delta_{n,m}\delta_{R,0}$.
\end{enumerate}

\textbf{Note:} Specific sets of two-dimensional lattices will be assigned by a random draw at the end of the class on 10/08/2026. Please make sure that the (two-student) project teams are established and that each team has (at least) one representative present for the draw.

\bigskip
\hrule
\bigskip

\section{General Considerations for the Solution}

To compute the band structure, we need to solve the Schr\"odinger equation in terms of the overlap integrals:
\[
	[\bar{\epsilon} \bar{\eta}(\vec{k}) + \bar{t}(\vec{k})]\bar{c} = E(\vec{k}) \bar{\eta}(\vec{k}) \bar{c}
\]
where the Hermitian matrices $\bar{\eta}$ and $\bar{t}$ are matrices with the following matrix elements:
\begin{align}
	\eta_{nm}(\vec{k}) &= \sum_{\vec{R}_{j}} \eta_{nm}(|\vec{R}_{j} + (\vec{\tau}_{m} - \vec{\tau}_{n})|) e^{i \vec{k} \cdot \vec{R}_{j}} \\
	t_{nm}(\vec{k}) &= \sum_{\vec{R}_{j}} t_{nm}(|\vec{R}_{j} + (\vec{\tau}_{m} - \vec{\tau}_{n})|) e^{i \vec{k} \cdot \vec{R}_{j}}
\end{align}
where $\vec{R}_{j}$ is a vector on the Bravais lattice, and $\vec{\tau}_{n}, \vec{\tau}_{m}$ are the sub-lattice basis vectors. 
$\bar{\epsilon}$ is the diagonal matrix of on-site energies for each orbital.

We can define $\bar{c}' = \bar{\eta}(\vec{k}) \bar{c}$ (such that $\bar{c} = \bar{\eta}^{-1}(\vec{k}) \bar{c}'$) and transform the equation into a standard eigenvalue problem:
\[
	[\bar{\epsilon} + \bar{t}(\vec{k}) \bar{\eta}^{-1}(\vec{k})] \bar{c}' = E(\vec{k}) \bar{c}'
\]
The procedure for both crystal A and crystal B is to determine the relevant $\bar{\epsilon}$, $\bar{\eta}(\vec{k})$, and $\bar{t}(\vec{k})$, compute the Fourier transformed expressions, and then solve the eigenvalue problem as a function of $\vec{k}$ to extract the band structure.

\bigskip
\hrule
\bigskip

\section{Crystal A: Square Lattice}

The square lattice has no lattice basis, therefore the only relevant $\vec{\tau}$ is $\vec{\tau} = (0,0)$, and there are only 2 orbitals. Therefore, the eigenvalue problem as a function of $\vec{k}$ is $2 \times 2$, corresponding to the two orbitals on-site.

Therefore, $\bar{\epsilon}$ is given by:
\[
	\bar{\epsilon} = \begin{pmatrix}
		0 & 0 \\
		0 & 0.8t_{0}
	\end{pmatrix}
\]
with $\bar{t}$ and $\bar{\eta}$ defined as $2 \times 2$ Hermitian matrices:
\[
	\bar{t}(\vec{k}) = \begin{pmatrix}
	t_{11}(\vec{k}) & t_{12}(\vec{k}) \\
	t_{12}^*(\vec{k}) & t_{22}(\vec{k})
	\end{pmatrix}, \qquad
	\bar{\eta}(\vec{k}) = \begin{pmatrix}
	\eta_{11}(\vec{k})  & \eta_{12}(\vec{k}) \\
	\eta_{12}^*(\vec{k}) & \eta_{22}(\vec{k})
	\end{pmatrix}
\]

\subsection*{On-Site Contribution ($R = 0$)}

The on-site contribution gives the identity matrix:
\[
	\eta_{nm}^{(0)}(\vec{k}) = \eta_{nm}(0) = \delta_{nm}, \qquad t_{nm}^{(0)}(\vec{k}) = 0
\]

\subsection*{Nearest Neighbor Contributions}

For the square lattice, this is straightforward. We consider one orbital pair as an example, and the rest follows by the same logic. 

Using the set of nearest-neighbor vectors $\vec{R}_{j} = (\pm a, 0), (0, \pm a)$ with distance $R = a$, since $\eta_{nm}$ and $t_{nm}$ depend only on distance $R$:
\begin{align*}
	\eta^{(1)}_{nm}(\vec{k}) &= \sum_{\vec{R}_{\text{NN}}} \eta_{nm}(\vec{R}_{j}) e^{i \vec{k} \cdot \vec{R}_{j}} = \eta_{nm}(a) \left[ e^{i k_{x}a} + e^{-i k_{x}a} + e^{i k_{y}a} + e^{-i k_{y}a} \right] \\
	&= 2 \eta_{nm}(a) \left[\cos(k_{x}a) + \cos(k_{y}a)\right] = 2 \eta_{nm}(a) \epsilon_{0}(\vec{k})
\end{align*}
where $\epsilon_{0}(\vec{k}) = \cos(k_{x}a) + \cos(k_{y}a)$. This factor is the same across all orbital pairs, so in general we can write:
\begin{align*}
	\eta^{(1)}_{nm}(\vec{k}) &= 2 \eta_{nm}(a) \epsilon_{0}(\vec{k}) \\
	t^{(1)}_{nm}(\vec{k}) &= 2 t_{nm}(a) \epsilon_{0}(\vec{k})
\end{align*}

\subsection*{Next-Nearest Neighbor Contributions}

A similar logic follows for the next-nearest neighbors at $\vec{R}_{j} = (\pm a, \pm a)$ with distance $R = \sqrt{2}a$:
\begin{align*}
	\eta^{(2)}_{nm}(\vec{k}) &= \eta_{nm}(\sqrt{2}a) \left[e^{i a(k_{x} + k_{y})} + e^{i a(k_{x} - k_{y})} + e^{-i a(k_{x} - k_{y})} + e^{-i a(k_{x} + k_{y})}\right] \\
	&= 2 \eta_{nm}(\sqrt{2}a)\left[\cos(a k_{x} + a k_{y}) + \cos(a k_{x} - a k_{y})\right] \\
	&= 4 \eta_{nm}(\sqrt{2}a) \cos(k_{x}a) \cos(k_{y}a) = 2 \eta_{nm}(\sqrt{2}a) \epsilon_{1}(\vec{k})
\end{align*}
where $\epsilon_{1}(\vec{k}) = \cos(a(k_{x} + k_{y})) + \cos(a(k_{x} - k_{y})) = 2 \cos(k_{x}a) \cos(k_{y}a)$.

The exact same is true for $t_{nm}$, so including the on-site term $\delta_{nm}$ for the overlap, we can write the total contributions as:
\begin{align*}
	\eta_{nm}(\vec{k}) &= \delta_{nm} + 2 \eta_{nm}(a) \epsilon_{0}(\vec{k}) + 2 \eta_{nm}(\sqrt{2}a) \epsilon_{1}(\vec{k}) \\
	t_{nm}(\vec{k}) &= 2 t_{nm}(a) \epsilon_{0}(\vec{k}) + 2 t_{nm}(\sqrt{2}a) \epsilon_{1}(\vec{k}) \\[1em]
	\epsilon_{0}(\vec{k}) &= \cos(k_{x}a) + \cos(k_{y}a) \\
	\epsilon_{1}(\vec{k}) &= 2 \cos(k_{x}a) \cos(k_{y}a)
\end{align*}

\bigskip
\hrule
\bigskip

\section{Crystal B: Centered Rectangular Lattice}

Crystal B is modeled on a centered rectangular lattice with a two-site basis:
\begin{itemize}
    \item Site 1 at $\vec{\tau}_1 = (0, 0)$, which has two orbitals (orbital 1 and orbital 2)
    \item Site 2 at $\vec{\tau}_2 = \left(\frac{a}{2}, \frac{b}{2}\right)$, which has one orbital (orbital 1)
\end{itemize}

This gives a 3-orbital basis:
\begin{enumerate}
    \item Orbital 1 on Site 1 ($\epsilon_1 = 0$)
    \item Orbital 2 on Site 1 ($\epsilon_2 = 0.8t_0$)
    \item Orbital 1 on Site 2 ($\epsilon_3 = \epsilon_1 = 0$)
\end{enumerate}

Since State 1 is orbital 1 (Site 1), State 2 is orbital 2 (Site 1), and State 3 is orbital 1 (Site 2), the off-diagonal inter-sublattice couplings map directly to the radial functions:
\begin{align*}
\eta_{13}(R) &\equiv \eta_{11}(R), \quad &t_{13}(R) &\equiv t_{11}(R) \\
\eta_{23}(R) &\equiv \eta_{12}(R), \quad &t_{23}(R) &\equiv t_{12}(R)
\end{align*}

Therefore, the eigenvalue problem as a function of $\vec{k}$ is $3 \times 3$, and $\bar{\epsilon}$ is given by:
\[
	\bar{\epsilon} = \begin{pmatrix}
		0 & 0 & 0 \\
		0 & 0.8t_0 & 0 \\
		0 & 0 & 0
	\end{pmatrix}
\]
with $\bar{t}$ and $\bar{\eta}$ defined as $3 \times 3$ Hermitian matrices:
\[
	\bar{t}(\vec{k}) = \begin{pmatrix}
	t_{11}(\vec{k}) & t_{12}(\vec{k})  & t_{13}(\vec{k}) \\
	t_{12}^*(\vec{k}) & t_{22}(\vec{k}) & t_{23}(\vec{k}) \\
	t_{13}^*(\vec{k}) & t_{23}^*(\vec{k}) & t_{33}(\vec{k})
	\end{pmatrix}, \qquad
	\bar{\eta}(\vec{k}) = \begin{pmatrix}
	\eta_{11}(\vec{k})  & \eta_{12}(\vec{k})  & \eta_{13}(\vec{k}) \\
	\eta_{12}^*(\vec{k}) & \eta_{22}(\vec{k}) & \eta_{23}(\vec{k}) \\
	\eta_{13}^*(\vec{k})  & \eta_{23}^*(\vec{k}) & \eta_{33}(\vec{k})
	\end{pmatrix}
\]

For this lattice, contributions are classified by the actual physical interatomic distance $R = |\vec{R}_{j} + (\vec{\tau}_{m} - \vec{\tau}_{n})|$ into on-site, nearest-neighbor, and next-nearest-neighbor shells.

\subsection*{1. On-Site ($R = 0$) Contribution}

Because Site 1 at $\vec{\tau}_{1} = (0,0)$ and Site 2 at $\vec{\tau}_{2} = \left(\frac{a}{2}, \frac{b}{2}\right)$ are physically distinct atoms, there is no $R = 0$ inter-sublattice term ($\eta_{13}^{(0)} = \eta_{23}^{(0)} = 0$ and $t_{nm}^{(0)} = 0$). The on-site overlap matrix is purely diagonal:
\[
	\eta_{nm}^{(0)}(\vec{k}) = \eta_{nm}(0) = \delta_{nm}, \qquad t_{nm}^{(0)}(\vec{k}) = 0
\]
Specifically:
\[
\eta_{11}^{(0)} = \eta_{22}^{(0)} = \eta_{33}^{(0)} = 1, \qquad \eta_{12}^{(0)} = \eta_{13}^{(0)} = \eta_{23}^{(0)} = 0
\]

\subsection*{2. Nearest Neighbor Contributions `(1)'}

\subsubsection*{Intra-Sublattice Elements ($11, 22, 33$ and $12$)}

For the intra-sublattice terms ($\vec{\tau}_m - \vec{\tau}_n = \vec{0}$), the nearest neighbors on the Bravais lattice lie along the axes at $\vec{R}_j = (\pm a, 0)$ and $(0, \pm b)$ with distances $R = a$ and $R = b$, respectively.

For the diagonal elements:
\begin{align*}
	\eta^{(1)}_{nn}(\vec{k}) &= 2 \eta_{nn}(a) \cos(k_x a) + 2 \eta_{nn}(b) \cos(k_y b) \\
	t^{(1)}_{nn}(\vec{k}) &= 2 t_{nn}(a) \cos(k_x a) + 2 t_{nn}(b) \cos(k_y b)
\end{align*}
where for State 3 (orbital 1 on Site 2), $\eta^{(1)}_{33} = \eta^{(1)}_{11}$ and $t^{(1)}_{33} = t^{(1)}_{11}$.

For the off-diagonal intra-sublattice elements ($\eta_{12}, t_{12}$), since orbitals 1 and 2 reside on the same site, the on-site overlap is zero by orthogonality. The nearest-neighbor interaction yields:
\begin{align*}
	\eta^{(1)}_{12}(\vec{k}) &= 2 \eta_{12}(a) \cos(k_x a) + 2 \eta_{12}(b) \cos(k_y b) \\
	t^{(1)}_{12}(\vec{k}) &= 2 t_{12}(a) \cos(k_x a) + 2 t_{12}(b) \cos(k_y b)
\end{align*}

\subsubsection*{Inter-Sublattice Elements ($\eta_{13}, t_{13}$ and $\eta_{23}, t_{23}$)}

For the inter-sublattice couplings, $\vec{\tau}_{2} - \vec{\tau}_{1} = \left( \frac{a}{2}, \frac{b}{2} \right)$. The four nearest-neighbor Site 2 atoms surrounding Site 1 lie symmetrically at displacements $\vec{d} = \left(\pm \frac{a}{2}, \pm \frac{b}{2}\right)$, all at the exact same physical distance $R = \frac{1}{2}\sqrt{a^{2} + b^{2}}$. These correspond to the four Bravais lattice vectors:
\[
\vec{R}_{j} \in \left\{ (0,0), (-a,0), (0,-b), (-a,-b) \right\}
\]

\begin{align*}
\eta_{13}^{(1)}(\vec{k}) &= \eta_{13}\left(\frac{\sqrt{a^{2} + b^{2}}}{2} \right) \left[ 1 + e^{- i a k_{x}} + e^{- i b k_{y}} + e^{- i (a k_{x} + b k_{y})} \right] \\
&= \eta_{13}\left(\frac{\sqrt{a^{2} + b^{2}}}{2} \right) \left(1 + e^{- i a k_{x}}\right)\left(1 + e^{- i b k_{y}}\right) \\[1em]
t_{13}^{(1)}(\vec{k}) &= t_{13}\left(\frac{\sqrt{a^{2} + b^{2}}}{2} \right) \left(1 + e^{- i a k_{x}}\right)\left(1 + e^{- i b k_{y}}\right) \\[1em]
\eta_{23}^{(1)}(\vec{k}) &= \eta_{23}\left(\frac{\sqrt{a^{2} + b^{2}}}{2} \right) \left(1 + e^{- i a k_{x}}\right)\left(1 + e^{- i b k_{y}}\right) \\[1em]
t_{23}^{(1)}(\vec{k}) &= t_{23}\left(\frac{\sqrt{a^{2} + b^{2}}}{2} \right) \left(1 + e^{- i a k_{x}}\right)\left(1 + e^{- i b k_{y}}\right)
\end{align*}

\subsection*{3. Next-Nearest Neighbor Contributions `(2)'}

\subsubsection*{Intra-Sublattice Elements ($11, 22, 33$ and $12$)}

Displacements to the four corner unit cells lie at $\vec{R}_j = (\pm a, \pm b)$ with distance $R = \sqrt{a^2 + b^2}$:

\begin{align*}
\eta_{nn}^{(2)}(\vec{k}) &= 4 \eta_{nn}\left(\sqrt{a^{2} + b^{2}} \right) \cos{\left(a k_{x} \right)} \cos{\left(b k_{y} \right)} \\
t_{nn}^{(2)}(\vec{k}) &= 4 t_{nn}\left(\sqrt{a^{2} + b^{2}} \right) \cos{\left(a k_{x} \right)} \cos{\left(b k_{y} \right)} \\[1em]
\eta_{12}^{(2)}(\vec{k}) &= 4 \eta_{12}\left(\sqrt{a^{2} + b^{2}} \right) \cos{\left(a k_{x} \right)} \cos{\left(b k_{y} \right)} \\
t_{12}^{(2)}(\vec{k}) &= 4 t_{12}\left(\sqrt{a^{2} + b^{2}} \right) \cos{\left(a k_{x} \right)} \cos{\left(b k_{y} \right)}
\end{align*}

\subsubsection*{Inter-Sublattice Elements ($\eta_{13}, t_{13}$ and $\eta_{23}, t_{23}$)}

The second inter-sublattice shell consists of 8 neighbors formed by:
\begin{itemize}
    \item 4 neighbors extended along $x$ at displacements $\vec{d} = \left(\pm \frac{3a}{2}, \pm \frac{b}{2}\right)$ with distance $R = \frac{1}{2}\sqrt{9a^{2} + b^{2}}$, generated by $\vec{R}_{j} \in \{(a,0), (a,-b), (-2a,0), (-2a,-b)\}$.
    \item 4 neighbors extended along $y$ at displacements $\vec{d} = \left(\pm \frac{a}{2}, \pm \frac{3b}{2}\right)$ with distance $R = \frac{1}{2}\sqrt{a^{2} + 9b^{2}}$, generated by $\vec{R}_{j} \in \{(0,b), (-a,b), (0,-2b), (-a,-2b)\}$.
\end{itemize}

\begin{align*}
\eta_{13}^{(2)}(\vec{k}) &= \eta_{13}\left(\frac{\sqrt{9 a^{2} + b^{2}}}{2} \right) \left[ e^{i a k_{x}} + e^{i(a k_{x} - b k_{y})} + e^{-2i a k_{x}} + e^{-i(2a k_{x} + b k_{y})} \right] \\
&\quad + \eta_{13}\left(\frac{\sqrt{a^{2} + 9 b^{2}}}{2} \right) \left[ e^{i b k_{y}} + e^{i(-a k_{x} + b k_{y})} + e^{-2i b k_{y}} + e^{-i(a k_{x} + 2b k_{y})} \right] \\[1em]
t_{13}^{(2)}(\vec{k}) &= t_{13}\left(\frac{\sqrt{9 a^{2} + b^{2}}}{2} \right) \left[ e^{i a k_{x}} + e^{i(a k_{x} - b k_{y})} + e^{-2i a k_{x}} + e^{-i(2a k_{x} + b k_{y})} \right] \\
&\quad + t_{13}\left(\frac{\sqrt{a^{2} + 9 b^{2}}}{2} \right) \left[ e^{i b k_{y}} + e^{i(-a k_{x} + b k_{y})} + e^{-2i b k_{y}} + e^{-i(a k_{x} + 2b k_{y})} \right] \\[1em]
\eta_{23}^{(2)}(\vec{k}) &= \eta_{23}\left(\frac{\sqrt{9 a^{2} + b^{2}}}{2} \right) \left[ e^{i a k_{x}} + e^{i(a k_{x} - b k_{y})} + e^{-2i a k_{x}} + e^{-i(2a k_{x} + b k_{y})} \right] \\
&\quad + \eta_{23}\left(\frac{\sqrt{a^{2} + 9 b^{2}}}{2} \right) \left[ e^{i b k_{y}} + e^{i(-a k_{x} + b k_{y})} + e^{-2i b k_{y}} + e^{-i(a k_{x} + 2b k_{y})} \right] \\[1em]
t_{23}^{(2)}(\vec{k}) &= t_{23}\left(\frac{\sqrt{9 a^{2} + b^{2}}}{2} \right) \left[ e^{i a k_{x}} + e^{i(a k_{x} - b k_{y})} + e^{-2i a k_{x}} + e^{-i(2a k_{x} + b k_y)} \right] \\
&\quad + t_{23}\left(\frac{\sqrt{a^{2} + 9 b^{2}}}{2} \right) \left[ e^{i b k_{y}} + e^{i(-a k_{x} + b k_{y})} + e^{-2i b k_{y}} + e^{-i(a k_{x} + 2b k_{y})} \right]
\end{align*}

The lower-triangle elements are given by the Hermitian conjugate relations:
\[
t_{ji}(\vec{k}) = t_{ij}^*(\vec{k}), \qquad \eta_{ji}(\vec{k}) = \eta_{ij}^*(\vec{k})
\]

\subsection*{4. Total Matrix Assembly \& Approximations}

The total matrix elements are combined as:
\begin{align*}
	\eta_{nm}(\vec{k}) &= \eta_{nm}^{(0)}(\vec{k}) + \eta_{nm}^{(1)}(\vec{k}) + \eta_{nm}^{(2)}(\vec{k}) \\
	t_{nm}(\vec{k}) &=  t^{(0)}_{nm}(\vec{k}) + t_{nm}^{(1)}(\vec{k}) + t_{nm}^{(2)}(\vec{k})
\end{align*}

\end{document}
